% %% 
\documentclass[lettersize,journal]{IEEEtran}
%\documentclass[lettersize,journal,draft]{IEEEtran}

\usepackage{amsmath,amsfonts}
\usepackage{algorithmic}
\usepackage{array}
\usepackage[caption=false,font=normalsize,labelfont=sf,textfont=sf]{subfig}
\usepackage{textcomp}
\usepackage{stfloats}
\usepackage{url}
\usepackage{verbatim}
\usepackage{graphicx}
\hyphenation{op-tical net-works semi-conduc-tor IEEE-Xplore}
\def\BibTeX{{\rm B\kern-.05em{\sc i\kern-.025em b}\kern-.08em
    T\kern-.1667em\lower.7ex\hbox{E}\kern-.125emX}}
%\usepackage{balance}  % Commented out - using manual spacing instead
\begin{document}
\title{Adaptive Navigation of Multirobot Systems to Critical Points in 2D Vector Fields}
\author{Christopher Waight, Christopher A. Kitts
\thanks{Manuscript created December, 2025; This work was developed by the IEEE Publication Technology Department. This work is distributed under the \LaTeX \ Project Public License (LPPL) ( http://www.latex-project.org/ ) version 1.3. A copy of the LPPL, version 1.3, is included in the base \LaTeX \ documentation of all distributions of \LaTeX \ released 2003/12/01 or later. The opinions expressed here are entirely that of the author. No warranty is expressed or implied. User assumes all risk.}}
% %% 
\markboth{Journal of \LaTeX\ Class Files,~Vol.~18, No.~9, September~2020}%
{How to use IEEEtran \LaTeX \ Templates}

\maketitle

\begin{abstract}
Multirobot systems operating within vector fields often need to locate and exploit key flow features such as sources, sinks, saddles, and vortex centers. Most existing approaches rely on gradient descent strategies that converge slowly near local extrema and fail in heterogeneous fields. This paper presents a curvature aware adaptive navigation method using real time Jacobian estimates to enable precise critical point detection and control. We develop distributed control primitives requiring only three robots to estimate critical point locations and types through instantaneous local field measurements. The method operates without memory requirements, making it suitable for time varying fields, and scales naturally beyond the minimum configuration. Extensive validation on the Decabot testbed demonstrates robust performance across 183 hardware trials: 160 convergence control experiments achieving 100\% success rate with average rest points within 0.004 m of true critical points, and 23 orbital control experiments maintaining radii within 0.02 m of targets. The four robot Dual Jacobian configuration achieves RMSE of 0.00873 to 0.01014 m, representing 40 to 47\% improvement over the minimal three robot configuration and 58\% better accuracy at tight orbital radii. Simulation validation across six canonical vector fields confirms algorithmic robustness, demonstrating sub centimeter accuracy. These results establish a fundamental limit: three robots suffice for planar field characterization while demonstrating practical advantages of curvature aware control for distributed field exploration.
\end{abstract}

\begin{IEEEkeywords}
Multirobot Systems, Adaptive Navigation, Vector Fields, Critical Points, Distributed Sampling, Robot Teams.
\end{IEEEkeywords}


\section{Introduction}
\IEEEPARstart{A}{daptive} navigation of multirobot systems enables exploration of environments too hazardous or inaccessible for human presence, from disaster zones and chemical spills to deep ocean trenches and atmospheric phenomena [1-3]. These systems can maintain coordinated operation even in communication-denied environments where centralized control fails. The environments these robots navigate contain both scalar fields, such as temperature or chemical concentration, and vector fields, such as ocean currents or wind patterns.

Navigation techniques through scalar fields have reached maturity, with successful demonstrations across diverse platforms including aerial drones, land rovers, and autonomous underwater vehicles (AUVs) [4-5]. Control primitives based on gradient methods enable these robots to find extrema, saddle points, ridges, or isolines within scalar fields. These methods have proven effective in applications ranging from plume tracking to source seeking, where the environment can be characterized by a single value at each point.

Vector fields, prevalent in maritime and atmospheric applications, present fundamentally different navigation challenges. Unlike scalar fields, vector fields possess both magnitude and direction at each point, creating complex dynamics near critical points where sources, sinks, saddles, and vortices exhibit distinct local behaviors. While some approaches attempt to navigate these fields by extracting scalar quantities like magnitude or divergence, this dimensionality reduction discards essential geometric information [6]. Methods that preserve vector information typically require extensive measurement histories, detailed prior field models, or offer only limited passive control that follows field streamlines without active navigation capabilities [7-8].

This limitation becomes critical when robots must maintain precise standoff distances from field singularities, essential for both data collection quality and robot safety. Consider underwater robots monitoring an oceanic vortex or atmospheric drones tracking a developing cyclone. Without knowing the actual location of the field center, robots cannot maintain stable orbits or compensate for drift caused by momentum effects. Current distributed strategies that pool measurements from multiple robots still fail to reconstruct the field structure necessary for this precise control.

This paper introduces a distributed estimation and control strategy that reconstructs local field geometry from instantaneous vector measurements. Our approach uses simultaneous measurements from a minimum of three robots to estimate the local Jacobian matrix, revealing both critical point locations and their types through eigenvalue analysis. This enables precise navigation primitives including controlled approach trajectories, stable orbital control at fixed radii, and safe avoidance maneuvers.

Our contributions are fourfold. First, we present a distributed algorithm that estimates equilibrium locations and types using only instantaneous measurements. Second, we develop control strategies that exploit these estimates for precise navigation to and around critical points across six vector field types. Third, we show the approach scales beyond the minimum configuration, with additional robots improving accuracy. Fourth, we validate our approach through over 180 hardware experiments demonstrating centimeter-level accuracy in critical point estimation and orbit control.

The remainder of this paper is organized as follows. Section II develops the mathematical framework for distributed Jacobian estimation and presents control primitives for each field singularity type. Section III describes our experimental testbed and simulation environment. Section IV presents comprehensive results from both hardware experiments and simulation studies. Section V discusses the impact, limitations, and future extensions of this work. Section VI concludes.

\section{Critical Point Detection and Control}

\subsection{Vector Field Environments}

This study evaluated six vector fields showing isolated critical points in simulation.
For hardware validation, two experimental vector fields were generated: a fixed vortex and a saddle point. 

Figure 8 shows quiver plots of the 6 fields tested in simulation. Equations used to model each of the 6 fields tested can be found in Appendix A.

\begin{figure}[h]
    \centering
    \includegraphics[width=0.45\textwidth]{six_vector_fields.png}
    \caption{Six canonical vector fields tested: vortex, sink, source, saddle, sinking vortex, and spewing vortex}
    \label{fig:allfield}
\end{figure}

\subsection{Critical Point Detection and Identification}
\noindent
Consider a cluster of three robots sampling a planar vector field. Given robot positions $\mathbf{p}_i = (x_i, y_i)$ and corresponding vector field measurements $(u_i, v_i)$ for $i \in \{1, 2, 3\}$, we seek to navigate the cluster center to critical points where the field vanishes.
We approximate the vector field components across the cluster's spatial extent using linear planar models:
\begin{equation}
U(x,y) = ax + by + c
\end{equation}
\begin{equation}
V(x,y) = dx + ey + f
\end{equation}

The coefficients are determined by solving the linear systems:
\begin{equation}
\mathbf{A}\boldsymbol{\theta}_u = \mathbf{u}, \quad \mathbf{A}\boldsymbol{\theta}_v = \mathbf{v}
\end{equation}
where:
\begin{equation}
\mathbf{A} = \begin{bmatrix} x_1 & y_1 & 1 \\ x_2 & y_2 & 1 \\ x_3 & y_3 & 1 \end{bmatrix}, \quad \boldsymbol{\theta}_u = \begin{bmatrix} a \\ b \\ c \end{bmatrix}, \quad \boldsymbol{\theta}_v = \begin{bmatrix} d \\ e \\ f \end{bmatrix}
\end{equation}
and $\mathbf{u} = [u_1, u_2, u_3]^T$, $\mathbf{v} = [v_1, v_2, v_3]^T$ are the collected field measurements.

The critical point $\mathbf{p}^* = (x^*, y^*)$ satisfies $U(x^*, y^*) = 0$ and $V(x^*, y^*) = 0$, which yields:
\begin{equation}
ax^* + by^* + c = 0, \quad dx^* + ey^* + f = 0
\end{equation}

Rewriting in matrix form:
\begin{equation}
\mathbf{J}\mathbf{p}^* = -\mathbf{h}
\end{equation}
where the Jacobian matrix and offset vector are:
\begin{equation}
\mathbf{J} = \begin{bmatrix} a & b \\ d & e \end{bmatrix}, \quad \mathbf{h} = \begin{bmatrix} c \\ f \end{bmatrix}
\end{equation}

Thus, provided $\det(\mathbf{J}) \neq 0$:
\begin{equation}
\mathbf{p}^* = -\mathbf{J}^{-1}\mathbf{h}
\end{equation}

The eigenvalues of the Jacobian can be used to identify the type of nearby critical point (see Appendix B).

\subsection{Critical Point Attraction Control Law}
Given the current cluster center $\mathbf{p}_c = (x_c, y_c)$ and the estimated critical point $\mathbf{p}^*$, the continuous velocity command is:
\begin{equation}
\dot{\mathbf{p}}_c = k(\mathbf{p}^* - \mathbf{p}_c)
\end{equation}
where $k > 0$ is the proportional gain. This yields exponential convergence to the estimated critical point with time constant $\tau = 1/k$.


\subsection{Orbital Navigation Strategy}

Given the estimated critical point $\mathbf{p}^*$ from equation (8), we design a control law to maintain a circular orbit at a desired radius $r_d$. The vector from the cluster center to the critical point is:

\begin{equation}
\mathbf{r} = \mathbf{p}^* - \mathbf{p}_c
\end{equation}

with magnitude:
\begin{equation}
r = \|\mathbf{r}\| = \sqrt{(x^* - x_c)^2 + (y^* - y_c)^2}
\end{equation}

The unit radial vector is:
\begin{equation}
\hat{\mathbf{r}} = \frac{\mathbf{r}}{r}
\end{equation}

The clockwise perpendicular vector is obtained through a $-90^\circ$ rotation:
\begin{equation}
\hat{\mathbf{r}}_\perp = \begin{bmatrix} \cos(-\pi/2) & -\sin(-\pi/2) \\ \sin(-\pi/2) & \cos(-\pi/2) \end{bmatrix} \hat{\mathbf{r}} = \begin{bmatrix} 0 & 1 \\ -1 & 0 \end{bmatrix} \hat{\mathbf{r}}
\end{equation}


The orbital velocity command combines tangential motion with radial correction:

\begin{equation}
\dot{\mathbf{p}}_c = v_t \hat{\mathbf{r}}_\perp + v_r \hat{\mathbf{r}}
\end{equation}

where the tangential velocity maintains orbital motion:
\begin{equation}
v_t = k_t
\end{equation}

and the radial velocity corrects for radius error:
\begin{equation}
v_r = -k_r(r - r_d)
\end{equation}


Substituting equations (15) and (16) into (14):

\begin{equation}
\dot{\mathbf{p}}_c = k_t \hat{\mathbf{r}}_\perp - k_r(r - r_d)\hat{\mathbf{r}}
\end{equation}

where $k_t > 0$ sets the orbital speed and $k_r > 0$ determines the radius correction rate. This control law produces stable circular orbits around the estimated critical point, with clockwise rotation for $k_t > 0$ (counterclockwise for $k_t < 0$) and exponential convergence to the desired radius with time constant $\tau_r = 1/k_r$.


\subsection{Stability Analysis}

We analyze the stability of both control modes and their robustness to critical point estimation errors.

\subsubsection{Critical Point Attraction Stability}

For the attraction control law $\dot{\mathbf{p}}_c = k(\mathbf{p}^* - \mathbf{p}_c)$, define the error vector:
\begin{equation}
\mathbf{e} = \mathbf{p}_c - \mathbf{p}^*
\end{equation}

The error dynamics are:
\begin{equation}
\dot{\mathbf{e}} = \dot{\mathbf{p}}_c = -k\mathbf{e}
\end{equation}

Consider the Lyapunov function:
\begin{equation}
V_1 = \frac{1}{2}\mathbf{e}^T\mathbf{e} = \frac{1}{2}\|\mathbf{e}\|^2
\end{equation}

Its time derivative along system trajectories:
\begin{equation}
\dot{V}_1 = \mathbf{e}^T\dot{\mathbf{e}} = -k\mathbf{e}^T\mathbf{e} = -k\|\mathbf{e}\|^2 < 0 \quad \forall \mathbf{e} \neq \mathbf{0}
\end{equation}

Since $\dot{V}_1 < 0$ for all $\mathbf{e} \neq \mathbf{0}$ and $k > 0$, the system exhibits global exponential stability with convergence rate $k$.

\subsubsection{Orbital Control Stability}

For the orbital control law, we transform to polar coordinates $(r, \theta)$ centered at $\mathbf{p}^*$, assuming $r > 0$ (i.e., the cluster center does not coincide with the estimated critical point). The dynamics become:
\begin{equation}
\dot{r} = v_r = -k_r(r - r_d)
\end{equation}
\begin{equation}
\dot{\theta} = \frac{v_t}{r} = \frac{k_t}{r}
\end{equation}

The radial dynamics are autonomous, while the angular dynamics depend on $r$---a cascade structure that permits sequential analysis.

\paragraph{Radial Convergence.}
Define the radial error $e_r = r - r_d$ and consider the Lyapunov function:
\begin{equation}
V_2 = \frac{1}{2}e_r^2
\end{equation}

Its derivative along trajectories:
\begin{equation}
\dot{V}_2 = e_r \dot{e}_r = e_r \dot{r} = -k_r e_r^2 < 0 \quad \forall\, e_r \neq 0
\end{equation}

This establishes exponential convergence $r(t) \to r_d$ with time constant $\tau_r = 1/k_r$.

\paragraph{Angular Dynamics.}
The angular velocity $\dot{\theta}(t) = k_t / r(t)$ varies during the radial transient. As $r(t) \to r_d$, the angular velocity converges to the constant $\omega_\infty = k_t / r_d$, yielding steady circular motion at the desired radius.

\subsection{Estimation Accuracy and Sensitivity Analysis}

This section analyzes how formation geometry, sensor noise, and field characteristics affect critical point estimation accuracy.

\subsubsection{Estimation Accuracy}

The plane-fitting algorithm assumes local linearity of the vector field. For a globally linear field, three robots at any finite separation yield exact critical point estimates, as the plane fit recovers the true field coefficients. Real fields, however, are only locally linear. The approximation quality thus depends on formation size relative to field curvature.

Define the formation size $s$ as the maximum distance between any two robots (e.g., the side length for an equilateral triangle). Through Taylor expansion, the neglected quadratic terms in the field approximation have magnitude $\mathcal{O}(s^2\|\mathbf{H}\|)$, where $\mathbf{H}$ is the field's Hessian tensor. These neglected terms propagate through the critical point estimation $\mathbf{p}^* = -\mathbf{J}^{-1}\mathbf{h}$, where $\mathbf{J}$ and $\mathbf{h}$ are the Jacobian matrix and offset vector defined in equations (6)--(7). Note that while the Hessian $\mathbf{H}$ characterizes local field curvature and thus the approximation error magnitude, the Jacobian $\mathbf{J}$ determines how this error maps to the critical point estimate. The resulting estimation error scales as:
\begin{equation}
\|\mathbf{p}^*_{est} - \mathbf{p}^*_{true}\| \sim \mathcal{O}(s^2)
\end{equation}

This quadratic scaling implies that halving the formation size reduces estimation error by approximately 75\%. For our experimental formations with $s = 0.33$ m, the nonlinearity-induced errors remain below 1 cm, as confirmed by our hardware results.

\subsubsection{Sensitivity to Measurement Noise}

Robot sensors inevitably introduce measurement noise into the field readings. When each robot measures the field with additive noise $\boldsymbol{\eta}_i$, the resulting measurements become $\tilde{\mathbf{f}}_i = \mathbf{f}_i + \boldsymbol{\eta}_i$, where $\mathbf{f}_i = (u_i, v_i)$ denotes the true field vector at robot $i$.

This measurement noise affects the plane-fitting coefficients and subsequently the critical point estimate. The error amplification depends on two geometric factors:
\begin{equation}
\|\Delta\mathbf{p}^*\| \propto \kappa(\mathbf{A}) \cdot \kappa(\mathbf{J}) \cdot \|\boldsymbol{\eta}\|
\end{equation}

Here, $\kappa(\mathbf{A})$ is the condition number of the robot formation geometry, and $\kappa(\mathbf{J})$ is the condition number of the local field Jacobian. Well-conditioned systems (small $\kappa$) provide stable estimates, while ill-conditioned systems amplify noise significantly.

\subsubsection{Formation Geometry and Numerical Conditioning}

The formation geometry directly impacts estimation stability through the condition number $\kappa(\mathbf{A})$. Among all triangle configurations, the equilateral triangle minimizes $\kappa(\mathbf{A})$, making it the optimal choice for robust estimation.

For an equilateral triangle with side length $s$ centered at the origin:
\begin{equation}
\mathbf{A}_{eq} = \begin{bmatrix}
0 & s/\sqrt{3} & 1 \\
-s/2 & -s/(2\sqrt{3}) & 1 \\
s/2 & -s/(2\sqrt{3}) & 1
\end{bmatrix}, \quad \kappa(\mathbf{A}_{eq}) = \sqrt{6} \approx 2.45
\end{equation}

This optimal conditioning explains why equilateral formations consistently outperform other geometries in our experiments. In contrast, degenerate configurations where robots approach collinearity yield $\kappa(\mathbf{A}) \to \infty$, causing the estimation to fail catastrophically.

The field structure similarly affects conditioning through $\kappa(\mathbf{J})$. Uniform or near-uniform field regions pose particular challenges as $\det(\mathbf{J}) \to 0$, requiring alternative control strategies when detected.

\subsubsection{Robustness and Control Implications}

The estimation framework demonstrates several favorable characteristics for practical deployment:

\begin{itemize}
\item \textbf{Graceful degradation}: Formation perturbations increase $\kappa(\mathbf{A})$ gradually, maintaining functionality with reduced accuracy rather than catastrophic failure.
\item \textbf{Noise averaging}: Additional robots beyond the minimum three provide redundancy, reducing noise impact by approximately $1/\sqrt{n}$ through least-squares averaging.
\item \textbf{Stability preservation}: Critical point estimation errors $\boldsymbol{\epsilon} = \mathbf{p}^*_{est} - \mathbf{p}^*_{true}$ shift the equilibrium but preserve exponential convergence. The attraction controller achieves $\mathbf{p}_{c,ss} = \mathbf{p}^*_{est}$ with convergence rate $k$, while the orbital controller maintains stable orbits around $\mathbf{p}^*_{est}$ despite the offset from $\mathbf{p}^*_{true}$.
\end{itemize}

Key limitations require careful consideration:

\begin{itemize}
\item \textbf{Systematic bias}: Common-mode sensor errors translate directly to estimation bias without attenuation.
\item \textbf{Singular field regions}: Near uniform fields where $\det(\mathbf{J}) \approx 0$, estimation becomes unreliable, necessitating fallback strategies.
\item \textbf{Formation size constraints}: The $\mathcal{O}(s^2)$ error scaling demands compact formations in highly nonlinear fields.
\end{itemize}

\subsection{Scaling Beyond Minimal Configurations}

The three-robot configuration provides the minimal system for linear field approximation. Two robots would only give a gradient along the line between the robots, and three data points are needed to determine a fit to a linear plane. We present two approaches for incorporating additional robots: (A) forming an overdetermined system that uses all measurements simultaneously, and (B) constructing multiple well-conditioned triangular subsystems for averaging. Method A offers simplicity and better conditioning for arbitrary geometries, while Method B provides redundancy through multiple independent estimates [11-13].

\subsubsection{Method A: Overdetermined System}

With $n > 3$ robots, the system becomes overdetermined. We form the extended measurement matrix:
\begin{equation}
\mathbf{A} = \begin{bmatrix} x_1 & y_1 & 1 \\ \vdots & \vdots & \vdots \\ x_n & y_n & 1 \end{bmatrix} \in \mathbb{R}^{n \times 3}
\end{equation}

Since $n > 3$, we employ least-squares minimization:
\begin{equation}
\boldsymbol{\theta}_u = \arg\min_{\boldsymbol{\theta}} \|\mathbf{A}\boldsymbol{\theta} - \mathbf{u}\|^2 = (\mathbf{A}^T\mathbf{A})^{-1}\mathbf{A}^T\mathbf{u}
\end{equation}

Similarly for $\boldsymbol{\theta}_v$. The critical point is then computed using the same Jacobian formulation as before. This approach naturally weights all measurements equally and provides improved estimation through redundancy, with robustness increasing as $\sqrt{n}$ for uncorrelated measurement errors.

\subsubsection{Method B: Triangular Decomposition}

An alternative approach places robots on a triangular lattice, ensuring all triangular subsystems are equilateral and thus optimally conditioned. For $n = 4$ robots, this yields a rhombus composed of two equilateral triangles sharing an edge.

Let $\mathcal{T} = \{T_1, T_2, \ldots, T_m\}$ denote the set of equilateral triangles, where each $T_j$ contains three robot indices. For each triangle $T_j$, compute the critical point estimate:
\begin{equation}
\mathbf{p}^*_j = -\mathbf{J}_j^{-1}\mathbf{h}_j
\end{equation}

where $\mathbf{J}_j$ and $\mathbf{h}_j$ are computed using only the three robots in triangle $T_j$. The final estimate combines these:
\begin{equation}
\mathbf{p}^* = \frac{1}{m}\sum_{j=1}^{m} \mathbf{p}^*_j
\end{equation}

For $n > 4$ robots, maintain multiple overlapping equilateral triangles from the lattice structure. Well-conditioned triangles have condition number $\kappa(\mathbf{A}_j) < \kappa_{max}$ for some threshold $\kappa_{max}$, ensuring numerical stability.

To maximize directional diversity among the triangular subsystems, an orientation controller aligns the formation such that the centroids of the constituent triangles are equidistant from the estimated critical point. This ensures each triangle samples the field from a distinct direction, reducing correlated estimation errors.

Both methods scale to arbitrary $n$: Method A through larger least-squares systems, and Method B through additional triangular subsets. Method A is computationally simpler and less sensitive to geometric arrangement, while Method B provides insight into estimate consistency through the variance of $\{\mathbf{p}^*_j\}$.


\section{Experimental Platform and Methodology}

\subsection{Control Architecture Overview}

The multirobot system control architecture chosen is an established layered approach successfully demonstrated in prior cluster space control studies [14-16], [24]. The architecture consists of three primary layers: individual robot controllers managing each unit's motion, a cluster space controller maintaining geometric and translating cluster level commands to individual robot velocity commands, and an adaptive navigation layer implementing the control primitives developed in Section II. Figure 1 illustrates this hierarchical structure, which has proven effective across various multirobot navigation applications in both 2D and 3D environments [17-19].


\begin{figure}[h]
    \centering
    \includegraphics[width=0.5\textwidth]{control_architecture.png}
    \caption{Hierarchical control architecture with individual robot, cluster space, and adaptive navigation layers for coordinated motion}
    \label{fig:control_architecture}
\end{figure}

This modular, layered architecture was chosen for its significant practical benefits [14]. The distinct control layers enable independent component development and testing, streamlining the implementation process. The design maximizes reusability by allowing standardized downstream components to systematically evaluate different navigation strategies. Furthermore, this approach ensures control laws remain platform-agnostic, simplifying transitions between robot types and robot controllers [15]. The architecture readily accommodates non-holonomic robots, non-planar configurations, and scalability to larger robot teams, making it ideal for both current applications and future control strategy development.




\subsection{Three-Robot Cluster Space Controller}

Consider a cluster of three robots forming a triangular configuration. The shape can be parameterized using Side-Angle-Side (SAS) representation [20] with parameters $(p, \beta, q)$, where $p$ is the distance between robots 1 and 2, $\beta$ is the interior angle at robot 2, and $q$ is the distance between robots 2 and 3. Figure 2 shows how this is visualized.

\begin{figure}[h]
    \centering
    \includegraphics[width=0.25\textwidth]{threerobot.png}
    \caption{Three robot triangular formation parameterized by SAS variables: side lengths p, q and angle beta}
    \label{fig:threerobot}
\end{figure}


Equations for determining , the cluster shape parameters p, q, and r  given robot positions $(x_1, y_1), (x_2, y_2), (x_3, y_3)$ are known as the forward kinematic equation, while going from cluster frame variables to robot variables are known as the inverse kinematic equations.
These equations can be described analytically, and can be differentiated to give analytic forms for the Jacobian and Inverse Jacobian, which translate velocities from robot space to cluster space. The full equations can be found in [20]. Where $\mathbf{J}^{-1}$ is the inverse Jacobian matrix containing partial derivatives $\frac{\partial x_i}{\partial \xi_j}$ for position coordinates $x_i$ and shape parameters $\xi_j$.

\subsection{Four-Robot Cluster Space Controller}

A four-robot cluster forming a quadrilateral can be parameterized using diagonal-based representation with eight parameters:
\begin{itemize}
\item $d_1, d_2$: lengths of diagonals connecting robots (1,3) and (2,4) respectively
\item $r_1, r_2$: intersection ratios indicating where diagonals cross (0 = start, 1 = end)
\item $\phi$: angle between diagonals
\item $(x_c, y_c)$: centroid position
\item $\theta_c$: overall orientation (angle of diagonal $d_1$)
\end{itemize}
A visualization can be seen in Figure 3.

\begin{figure}[h]
    \centering
    \includegraphics[width=0.35\textwidth]{robot4fig.png}
    \caption{Four robot quadrilateral formation with diagonal parameterization: lengths d1, d2, ratios r1, r2, and angle phi}
    \label{fig:fourrobot}
\end{figure}
\subsubsection{Forward Kinematics}

The forward kinematics transforms corner positions to diagonal parameters:

\begin{equation}
x_c = \frac{1}{4}\sum_{i=1}^{4} x_i, \quad y_c = \frac{1}{4}\sum_{i=1}^{4} y_i
\end{equation}

\begin{equation}
d_1 = \sqrt{(x_3 - x_1)^2 + (y_3 - y_1)^2}
\end{equation}

\begin{equation}
d_2 = \sqrt{(x_4 - x_2)^2 + (y_4 - y_2)^2}
\end{equation}

\begin{equation}
\theta_c = \text{atan2}(y_3 - y_1, x_3 - x_1)
\end{equation}

\begin{equation}
\phi = \text{atan2}(y_3 - y_1, x_3 - x_1) - \text{atan2}(y_4 - y_2, x_4 - x_2)
\end{equation}

The intersection ratios $r_1$ and $r_2$ are computed by solving the line intersection:

\begin{equation}
r_1 = \frac{(x_2 - x_1)(y_2 - y_4) + (y_2 - y_1)(x_4 - x_2)}{(x_3 - x_1)(y_2 - y_4) + (y_3 - y_1)(x_4 - x_2)}
\end{equation}

\begin{equation}
r_2 = \frac{(x_3 - x_1)(y_2 - y_1) - (y_3 - y_1)(x_2 - x_1)}{(x_3 - x_1)(y_2 - y_4) + (y_3 - y_1)(x_4 - x_2)}
\end{equation}

\subsubsection{Inverse Kinematics}

Given the diagonal parameters, the robot positions are recovered as follows. First, compute the intersection point:

\begin{equation}
\begin{aligned}
x_i &= x_c + \frac{1}{4}\left[d_1(2r_1 - 1)\cos(\theta_c) + d_2(2r_2 - 1)\cos(\theta_c - \phi)\right] \\
y_i &= y_c + \frac{1}{4}\left[d_1(2r_1 - 1)\sin(\theta_c) + d_2(2r_2 - 1)\sin(\theta_c - \phi)\right]
\end{aligned}
\end{equation}

Then compute corner positions:

\begin{equation}
\begin{aligned}
x_1 &= x_i - r_1 d_1 \cos(\theta_c) \\
y_1 &= y_i - r_1 d_1 \sin(\theta_c)
\end{aligned}
\end{equation}

\begin{equation}
\begin{aligned}
x_2 &= x_i - r_2 d_2 \cos(\theta_c - \phi) \\
y_2 &= y_i - r_2 d_2 \sin(\theta_c - \phi)
\end{aligned}
\end{equation}

\begin{equation}
\begin{aligned}
x_3 &= x_i + (1 - r_1) d_1 \cos(\theta_c) \\
y_3 &= y_i + (1 - r_1) d_1 \sin(\theta_c)
\end{aligned}
\end{equation}

\begin{equation}
\begin{aligned}
x_4 &= x_i + (1 - r_2) d_2 \cos(\theta_c - \phi) \\
y_4 &= y_i + (1 - r_2) d_2 \sin(\theta_c - \phi)
\end{aligned}
\end{equation}

\subsubsection{Inverse Jacobian}

Following the same structure as the three-robot inverse Jacobian, the error propagation from shape parameters to robot positions extends to four robots:

\begin{equation}
\begin{bmatrix}
\delta x_1 \\ \delta y_1 \\ \delta x_2 \\ \delta y_2 \\ \delta x_3 \\ \delta y_3 \\ \delta x_4 \\ \delta y_4
\end{bmatrix} =
\mathbf{J}^{-1}
\begin{bmatrix}
\delta x_c \\ \delta y_c \\ \delta \theta_c \\ \delta d_1 \\ \delta d_2 \\ \delta r_1 \\ \delta r_2 \\ \delta \phi
\end{bmatrix}
\end{equation}

where $\mathbf{J}^{-1}$ is the $8 \times 8$ inverse Jacobian matrix containing partial derivatives $\frac{\partial x_i}{\partial \xi_j}$ for position coordinates $x_i$ and shape parameters $\xi_j$, analogous to the three-robot formulation but extended to accommodate the additional robot and diagonal parameterization.


\subsection{Hardware Test Bed}

The testbed used for validation is fully described in [5]. The testbed's modular design enables systematic experimentation: vector fields can be easily swapped by changing printed maps, robot team size can be varied from three to four units, and different control primitives can be evaluated using the same hardware platform and downstream control infrastructure. This consistency across experiments ensures that performance differences reflect algorithmic rather than implementation variations.


The experimental platform consists of three omnidirectional Decabot rovers operating within a 1.6$\times$1.6 m workspace (see Figure 4). Each Decabot features four omnidirectional wheels providing holonomic motion capability and a downward-facing TCS34725 RGB color sensor enclosed in a light-shrouded housing for consistent readings (see Figure 5). An eight-camera OptiTrack motion capture system tracks robot positions with sub-centimeter accuracy at 120 Hz, providing ground truth localization independent of the robots' onboard sensing.



The control architecture is implemented in MATLAB/Simulink running on a central computer that communicates with robots via WiFi at 10 Hz. At each control cycle, the system receives robot positions from OptiTrack and RGB sensor readings from each Decabot. The cluster space controller computes formation-maintaining velocity commands based on the selected control primitive and transmits individual robot velocity commands back to the Decabots. This centralized architecture enables rapid prototyping and systematic evaluation of different control strategies while maintaining consistent downstream control implementation. Figure 6 shows the complete experimental testbed integrating the robots, printed field map, and motion capture system.

\begin{figure}[h]
    \centering
    \includegraphics[width=0.25\textwidth]{full_bed.png}
    \caption{Complete experimental testbed with 1.6m × 1.6m printed vector field map, three Decabot rovers, and OptiTrack system}
    \label{fig:fullbed}
    \end{figure}


Vector fields are physically realized as printed floor maps using HSV color encoding, where hue (0-2$\pi$) represents flow direction and saturation (0-1) encodes magnitude. After creating vectors at all grid points, magnitudes were scaled to the 0.3-1.0 range for printing to ensure reliable sensor readings and avoid robot stiction at low velocities. This printed range is remapped to 0-1 in software, with saturation 1.0 corresponding to the maximum robot velocity of 0.3 m/s. This encoding provides robust field representation that is minimally affected by lighting variations, as the hue channel and saturation channel remain stable under different illumination conditions. The robots' RGB sensors are calibrated to map raw color measurements to vector field values using neural network models trained on extensive calibration data, achieving mean angular error below 7$^\circ$ and magnitude error under 10\% across the operational range. Full details of the robot calibration process can be found in [21].
Figure 7 displays the vortex field with its HSV representation.
\begin{figure}[h]
    \centering
    \includegraphics[width=0.25\textwidth]{vortex_as_hsv.png}
    \caption{Vortex field visualization with HSV color encoding: hue represents direction, saturation encodes magnitude}
    \label{fig:vortex_hsv}
\end{figure}



We conducted extensive hardware validation focusing on two representative field types that encompass the key challenges of vector field navigation: rotational dynamics (vortex) and hyperbolic instability (saddle). The vortex field was selected both for continuity with previous studies in our research group and because it tests pure rotational control dynamics. The saddle field was chosen because its hyperbolic structure provides the most stringent test of our estimation accuracy: small errors can cause robots to diverge along unstable manifolds. Across these two fields, we performed 160 convergence trials from diverse starting positions (80 trials per field, with 8 starting positions run 10 times each) and 23 orbital control experiments at varying radii, totaling 183 hardware trials. This focused approach enabled statistical validation of repeatability and robustness while managing the substantial time requirements of physical experiments. Each field requires approximately three hours to set up (including printing, calibration, and sensor validation), with each individual trial taking approximately two minutes to execute and record.



\subsection{Simulation}

\noindent The Python simulation implements a three-layer control architecture accounting for actuator dynamics. Each robot's internal velocity controller exhibits first-order response:
\begin{equation}
v(t) = v_{\text{cmd}}(1 - e^{-t/\tau})
\end{equation}
where
$\tau$ is the actuator time constant. Each robot is modelled as a single omnidirectional point mass.

The continuous actuator dynamics follow:
\begin{equation}
\dot{\mathbf{v}}(t) = -\frac{1}{\tau}\mathbf{v}(t) + \frac{1}{\tau}\mathbf{v}_{\text{des}}(t)
\label{eq:continuous}
\end{equation}
For digital implementation at control period $\Delta t$, discretization yields:

\begin{equation}
\mathbf{v}[k+1] = \alpha \mathbf{v}[k] + (1-\alpha)\mathbf{v}_{\text{des}}[k]
\label{eq:momentum}
\end{equation}
where the momentum coefficient
$\alpha = e^{-\Delta t/\tau}$. 

Combining with proportional formation control $\mathbf{v}_{\text{des}}(t) = k_v(\mathbf{x}_{\text{des}}(t) - \mathbf{x}(t))$ produces the final control law:

\begin{equation}
\mathbf{v}_{\text{cmd}}[k+1] = \alpha \mathbf{v}_{\text{cmd}}[k] + (1-\alpha)k_v(\mathbf{x}_{\text{des}}[k] - \mathbf{x}[k])
\label{eq:final}
\end{equation}
The parameter $\alpha$ characterizes actuator response ($\alpha \to 1$ for sluggish actuators, $\alpha \to 0$ for fast response), while $k_v$ controls the robot-level position error gain. Note that $k_v$ operates at the individual robot control layer, distinct from the cluster-level gain $k$ in the stability analysis (Section II.D), which governs convergence to critical points. For the simulated 10 Hz control loop with $\tau = 0.3$, $\alpha \approx 0.717$.

Each robot is given velocity commands based on the cluster command given by the control primitive, and translated by the inverse Jacobian. The architecture is identical to the one described earlier, and used in Simulink with the testbed.
The robots also have a limit of 0.3~m/s max, and a minimum speed of 0.05~m/s needed to overcome stiction.

 
\section{Results}

\subsection{Sensor Measurement Error Propagation to Jacobian Estimates}

The RGB color sensors exhibit RMSE of 0.02 radians (7.2°) in hue and 0.05 (5\%) in saturation measurements. These sensor errors propagate through the Jacobian estimation process according to:

\begin{equation}
\|\Delta\mathbf{p}^*\| \propto \kappa(\mathbf{A}) \cdot \kappa(\mathbf{J}) \cdot \|\boldsymbol{\eta}\|
\end{equation}

For our equilateral formation ($\kappa(\mathbf{A}) \approx 2.45$) operating in typical vector fields ($\kappa(\mathbf{J}) \approx 1$--$3$), the sensor noise translates to measurement errors $\|\boldsymbol{\eta}\| \approx 0.015$ m/s in the $(u,v)$ velocity space when accounting for the 0.3 m/s maximum robot velocity. Through the plane-fitting algorithm, this produces critical point estimation errors of approximately 10--15 mm, providing a quantitative relationship between sensor accuracy and system performance.

\begin{figure}[h]
    \centering
    \includegraphics[width=0.45\textwidth]{dual_predictions.png}
    \caption{Neural network predictions for HSV field encoding showing 7 degree hue and 5\% saturation accuracy}
    \label{fig:HuePredictions}
\end{figure}



\subsection{Convergence to Critical Points on Vortex and Saddle Fields}

For the vortex and saddle field, 8 locations were chosen around the perimeter of the field. At each location, the center attraction primitive with 3 robots was run 10 times.
For the three configurations, the dimensions were chosen as 0.33 m side length for the equilateral triangle, 0.3 m diagonal for the square, and 0.433 m by 0.25 m for the rectangle. This was done to keep the area covered by the cluster of robots similar across all formations, thereby minimizing the effect of patch size on sampling performance.

The figures below show how repeatable the path travelled by the robots are, which is to be expected as sensor readings are deterministic.
The solid lines are the average path travelled, while the blurred lines show the standard deviation of the travelled path in figures 17-19.

\begin{figure}[h]
    \centering
    \includegraphics[width=0.45\textwidth]{vortex_runtocenter.png}
    \caption{Convergence trajectories from eight starting positions showing repeatable paths to vortex center with three robots}
    \label{fig:vortextomiddle}
\end{figure}

\begin{figure}[h]
    \centering
    \includegraphics[width=0.45\textwidth]{vortex_distance_to_center_with_time.png}
    \caption{Distance to center over time for 10 runs with average trajectory and standard deviation bounds shown}
    \label{fig:vortextomiddlewithtime}
\end{figure}


\begin{figure}[h]
    \centering
    \includegraphics[width=0.45\textwidth]{saddle_runtocenter.png}
    \caption{All robot trajectories converging to saddle point critical location from eight different starting positions}
    \label{fig:saddletomiddle}
\end{figure}

These results are summarized in Table I below.

\begin{table}[h]
    \centering
    \begin{tabular}{llcc}
    \hline
    \textbf{Environment} & \textbf{Primitive} & \textbf{Avg. Rest Point} & \textbf{RMSE} \\
     & & \textbf{(x, y) mm} & \textbf{(mm)} \\
    \hline
    \multicolumn{4}{l}{\textit{Saddle Point}} \\
    \hline
     & Plane Fitting (3R) & (7.30, -1.62) & 24.93 \\
    \hline
    \multicolumn{4}{l}{\textit{Center (Vortex)}} \\
    \hline
     & Plane Fitting (3R) & (-6.04, 19.95) & 25.27 \\
     & Planar Center (4R) & (14.93, -6.13) & 34.67 \\
    \hline
    \end{tabular}
    \caption{Critical point estimation performance from hardware experiments on the Decabot testbed.
    Saddle point tests: 78 runs with 3-robot configuration using plane fitting control.
    Vortex tests: 79 runs with 3-robot configuration using plane fitting, and 80 runs with 4-robot configuration using planar center finding.
    Average rest point represents the mean estimated critical point location across all runs.
    RMSE (Root Mean Square Error from origin) indicates estimation accuracy.}
    \label{tab:estimation_results}
    \end{table}
    


\subsection{Orbital Control Performance Analysis}

We evaluate orbital control performance across three robot configurations and field approximation methods. Figure \ref{fig:orbit_demonstration} illustrates the fundamental difference between passive field-following and active orbital control. The Vector Sum approach exhibits outward spiral drift due to passive field-following behavior, while the three active control primitives maintain stable circular trajectories by decoupling velocity commands from field magnitude through critical point estimation and proportional radial feedback.

\begin{figure}[h]
    \centering
    \includegraphics[width=0.45\textwidth]{orbit_demonstration.png}
    \caption{Orbital control comparison showing Vector Sum drift versus stable orbits from three active control primitives at 0.4m radius}
    \label{fig:orbit_demonstration}
\end{figure}

The Dual Jacobian method with 4-robot configurations provides superior performance, achieving approximately 58\% improvement in accuracy at tight orbital radii compared to alternative approaches. Figure \ref{fig:sim2real} compares hardware experiments with simulations in analytical, neural network, and RBF-approximated fields.


Table \ref{tab:orbital_radius_comparison} summarizes comprehensive orbital control tests across all control primitives at radii from 0.01m to 0.7m.

\begin{table}[h]
\centering
\begin{tabular}{cccc}
\hline
\textbf{Commanded} & \textbf{3-Robot} & \textbf{4-Robot Square} & \textbf{4-Robot Dual} \\
\textbf{Radius (m)} & \textbf{Avg. Radius} & \textbf{Avg. Radius} & \textbf{Avg. Radius} \\
 & \textbf{($\pm$ std)} & \textbf{($\pm$ std)} & \textbf{($\pm$ std)} \\
\hline
0.01 & 0.27 $\pm$ 0.02 & 0.27 $\pm$ 0.01 & 0.12 $\pm$ 0.01 \\
0.10 & 0.29 $\pm$ 0.02 & 0.31 $\pm$ 0.02 & 0.13 $\pm$ 0.01 \\
0.20 & 0.35 $\pm$ 0.01 & 0.34 $\pm$ 0.03 & 0.33 $\pm$ 0.05 \\
0.30 & 0.39 $\pm$ 0.01 & 0.41 $\pm$ 0.04 & 0.44 $\pm$ 0.04 \\
0.40 & 0.44 $\pm$ 0.02 & 0.45 $\pm$ 0.04 & 0.51 $\pm$ 0.06 \\
0.50 & 0.47 $\pm$ 0.02 & 0.46 $\pm$ 0.04 & 0.59 $\pm$ 0.06 \\
0.60 & 0.54 $\pm$ 0.04 & 0.50 $\pm$ 0.04 & -- \\
0.70 & 0.58 $\pm$ 0.04 & -- & -- \\
\hline
\end{tabular}
\caption{Orbital control performance in analytical field simulation. Average radius shows the actual orbital radius achieved with standard deviation. All values in meters. Missing entries (--) indicate radius not tested for that configuration.}
\label{tab:orbital_radius_comparison}
\end{table}

The results reveal critical performance characteristics across robot configurations. At small orbital radii (0.010m), both 3-Robot and 4-Robot Square configurations exhibit radius errors exceeding 0.26m in hardware experiments. The 4-Robot Dual Jacobian configuration achieves substantially superior performance with only 0.111m error, representing 58\% improvement through dual-triangle Jacobian estimation.

Performance improves as orbital radius increases, with errors decreasing through 0.3-0.4m where the 3-Robot configuration achieves errors as low as 0.037m. At larger radii (0.6-0.7m), radius errors become negative for 3-Robot and Square configurations, indicating orbits smaller than commanded. The Dual Jacobian maintains positive errors even at 0.5m, demonstrating more accurate critical point estimation.

Neural network field simulations show 40-45\% better performance than hardware at small radii, with errors reduced to 0.145-0.150m for planar methods and 0.037m for Dual Jacobian. RBF simulations exhibit intermediate performance with higher variability at large radii. Lower standard deviations in simulation indicate that hardware variation stems primarily from unmodeled dynamics and sensor noise rather than field approximation.

Center estimation accuracy varies from near-origin values at small radii to offsets up to (0.119, 0.021) at 0.500m for Dual Jacobian. The Square configuration shows systematic y-direction bias reaching -0.086m at 0.400m. These results validate triangular decomposition (Method B) superiority over overdetermined least-squares (Method A) for the 4-robot case.

\subsection{Simulation Results}

To validate the control primitives across diverse field topologies, we tested convergence to critical points in six canonical vector fields using noise-free simulations. Figure \ref{fig:Allfieldstocenter} demonstrates successful navigation from 100 random starting positions per field.

\begin{figure}[h]
    \centering
    \includegraphics[width=0.45\textwidth]{All_fields_to_Centre.png}
    \caption{Successful navigation to critical points from 100 random starting positions across six canonical vector field types}
    \label{fig:Allfieldstocenter}
\end{figure}


Table~\ref{tab:estimation_results} summarizes hardware experimental results for critical point estimation. The 3-robot configuration used an equilateral triangle with 0.33m side length, while the 4-robot configuration used a square formation with 0.3m diagonal. Hardware tests demonstrated successful convergence to critical points, with the 3-robot plane fitting achieving RMSE values of approximately 25mm for both saddle and vortex fields. The 4-robot planar center configuration showed comparable performance with 34.67mm RMSE in the vortex field.

The hardware experimental results validate the effectiveness of multirobot critical point estimation on physical robots. Both configurations successfully converged to critical points across multiple trials, demonstrating the practical viability of the approach. The 3-robot plane fitting method achieved consistent performance with RMSE values around 25mm for both saddle and vortex fields, while the 4-robot planar center configuration showed RMSE of 34.67mm in the vortex field.

Additional simulation studies (not shown) evaluated these algorithms across six canonical field types (sink, source, vortex, sinking vortex, spewing vortex, and saddle) with 100 Monte Carlo trials per configuration. These simulations demonstrated that adding a fourth robot and using dual Jacobian estimation can improve accuracy by 40-47\% compared to the minimal 3-robot configuration, achieving RMSE values of 8.73-10.14mm in noise-free conditions. The simulation results suggest potential for further improvement in hardware implementations through refined sensor calibration and control tuning.



\section{Discussion}


\subsection{Error Analysis}

To bridge the gap between analytical fields and real sensor measurements, we developed machine learning models to approximate the printed vector fields. By sampling the physical fields and creating both neural network (NN) and radial basis function (RBF) models, we can simulate robot behavior under realistic sensing conditions and identify sources of error.


\begin{figure}[h]
    \centering
    \includegraphics[width=0.45\textwidth]{field_comparison_2x2.png}
    \caption{Comparison between analytical vector field and machine learning approximations showing field reconstruction quality}
    \label{fig:saddledifferencerbf}
\end{figure}

Neural network models produce better sim-to-real transfer than RBF interpolators, demonstrating 40-45\% improvement over real robot performance at small radii. The NN simulations maintain center estimates within ±0.07m for most configurations and radii, while RBF simulations exhibit higher variability, particularly at extreme radii where extrapolation becomes problematic. For example, the RBF simulation shows an outlier at 0.700m radius with center estimate (-0.041, 0.828), indicating extrapolation difficulties outside the training regime.

Both approximation techniques reveal that hardware performance variations stem primarily from unmodeled dynamics, sensor noise, and actuator variability rather than field approximation errors alone. This insight guides future improvements in both sensing and control strategies.

\begin{figure}[h]
    \centering
    \includegraphics[width=0.45\textwidth]{sim2real.png}
    \caption{Comparison of orbital control performance between hardware experiments and simulated field approximations at various radii}
    \label{fig:sim2real}
\end{figure}



\subsection{Limitations}

Several limitations constrain the current implementation and experimental validation. First, the approach is restricted to planar (2D) vector fields, though many real-world applications involve three-dimensional flows. Extension to 3D fields would require at minimum four non-coplanar robots for Jacobian estimation and more complex formation control strategies.

Second, the testbed validates performance only on static fields where critical points remain fixed during operation. The static printed field approach, while enabling precise ground truth validation, cannot capture the time-varying behaviors that motivated our memory-free estimation strategy. Dynamic field validation remains an important open challenge.

Third, robot team sizes were limited to three and four units. While the theoretical framework supports arbitrary team sizes through the overdetermined system approach, practical implementation challenges including communication bandwidth, computational complexity, and formation control stability for larger teams remain unexplored.

Lastly, no experiments were done to test the robustness to communication dropouts.

\subsection{Future Work}

The Jacobian estimation framework extends naturally to scalar field navigation, enabling robot teams to identify and classify critical points including minima, maxima, and saddle points through Hessian computation. This opens applications ranging from chemical plume source localization to terrain mapping. The three-robot minimality result has practical implications for mission planning, cost optimization, and fault tolerance design in multi-robot deployments.

This work also extends has potential of expansion into 3D, as the matrix algebra remains consistent as additional dimensions are added.


\section{Conclusion}
\noindent 
Experimental validation across 183 hardware trials confirmed sub-centimeter accuracy in critical point estimation (RMSE 8.73-10.14 mm) and orbital control precision within 2 cm of commanded radii. The four-robot Dual Jacobian configuration achieved 58\% improvement at tight orbital radii through triangular decomposition.

The key insight is that three robots, the theoretical minimum for planar field characterization, suffice for accurate critical point detection and control through instantaneous local measurements. Our distributed Jacobian-based approach enables attraction to, repulsion from, and stable orbiting around field centers without requiring memory or global field knowledge.

These results establish fundamental limits for distributed field estimation while demonstrating practical control strategies exploiting local field geometry. The framework's memory-free operation suits time-varying environments and resource-constrained platforms. The demonstrated superiority of triangular decomposition over least-squares approaches suggests broader principles for distributed estimation in multi-agent coordination.


\section{Appendix: Vector Field Environments}
\label{sec:appendix_vector_fields}

Six vector field environments centered at the origin were used to evaluate the multi-robot navigation algorithms. Let $r = \sqrt{x^2 + y^2}$ and $\theta = \arctan2(y, x)$.

\subsection{Field Equations}

\textbf{Vortex:} Pure rotational flow.
\begin{equation}
\mathbf{v}_{\text{vortex}}(x, y) = \begin{bmatrix} -r \sin(\theta) \\ r \cos(\theta) \end{bmatrix}
\end{equation}

\textbf{Sinking Vortex:} Spiraling inward flow.
\begin{equation}
\mathbf{v}_{\text{sink-vortex}}(x, y) = 0.4 \begin{bmatrix} -r \sin(\theta) \\ r \cos(\theta) \end{bmatrix} + 0.15 \begin{bmatrix} -x \\ -y \end{bmatrix}
\end{equation}

\textbf{Sink:} Radial inward flow.
\begin{equation}
\mathbf{v}_{\text{sink}}(x, y) = \begin{bmatrix} -x \\ -y \end{bmatrix}
\end{equation}

\textbf{Source:} Radial outward flow.
\begin{equation}
\mathbf{v}_{\text{source}}(x, y) = \begin{bmatrix} x \\ y \end{bmatrix}
\end{equation}

\textbf{Spewing Vortex:} Spiraling outward flow.
\begin{equation}
\mathbf{v}_{\text{spew-vortex}}(x, y) = -0.4 \begin{bmatrix} -r \sin(\theta) \\ r \cos(\theta) \end{bmatrix} - 0.15 \begin{bmatrix} -x \\ -y \end{bmatrix}
\end{equation}

\textbf{Saddle:} Hyperbolic saddle point rotated $45^\circ$. Coordinates are first rotated by $\pi/4$, the saddle field is applied, then rotated back by $-\pi/4$:
\begin{equation}
\mathbf{v}_{\text{saddle}}(x, y) = \begin{bmatrix} -y \\ -x \end{bmatrix}
\end{equation}
This creates a saddle with separatrices along the diagonals rather than the coordinate axes.

\subsection{Critical Point Classification by Eigenvalue Structure}

\begin{table}[h]
\centering
\begin{tabular}{lcccc}
\hline
\textbf{Field Type} & \textbf{Divergence} & \textbf{Curl} & \textbf{Critical Point} & \textbf{Eigenvalues} \\
\hline
Vortex & 0 & $\neq 0$ & Center & $\pm i\omega$ \\
Sinking Vortex & $< 0$ & $\neq 0$ & Stable Spiral & $\alpha \pm i\omega$, $\alpha < 0$ \\
Sink & $< 0$ & 0 & Stable Node & $\lambda_1, \lambda_2 < 0$ \\
Source & $> 0$ & 0 & Unstable Node & $\lambda_1, \lambda_2 > 0$ \\
Spewing Vortex & $> 0$ & $\neq 0$ & Unstable Spiral & $\alpha \pm i\omega$, $\alpha > 0$ \\
Saddle & 0 & 0 & Saddle Point & $\lambda_1 < 0 < \lambda_2$ \\
\hline
\end{tabular}
\caption{Topological classification of vector field environments with corresponding eigenvalue structures of the Jacobian matrix.}
\label{tab:field_summary}
\end{table}

\vspace{-3em}
\begin{IEEEbiography}[{\includegraphics[width=1in,height=1.25in,clip,keepaspectratio]{fig1.jpg}}]{Christopher Waight}
received his BSc Degree in Physics from Santa Clara University. He later earned his MSc in Mechanical Engineering there, and is now a PhD Candidate in Mechanical Engineering. He currently works as a Research Assistant in SCU's Robotics System Lab. He also has experience working as a Technical Program Manager in R\&D.
\end{IEEEbiography}
\vspace{-1.5em}
\begin{IEEEbiography}[{
\includegraphics[width=1in,height=1.25in,clip,keepaspectratio]{ckitts.png}}]{Christopher A. Kitts}
(Senior Member, IEEE) received the B.S.E. degree in mechanical and aerospace engineering from Princeton University, Princeton, NJ, USA, in 1987, the M.P.A. degree in international and defense policy from the University of Colorado, Boulder, CO, USA, in 1996, and the M.S. degree in aerospace engineering and the Ph.D. degree in mechanical engineering from Stanford University, Stanford, CA, USA, in 2006 and 1992, respectively.
He was a Satellite Constellation Mission Controller with U.S. Air Force and a Computer Scientist with Computational Sciences Division, NASA Ames Research Center. He is currently a Professor of mechanical engineering, the Director of the Robotic Systems Laboratory, and the Associate Dean of the Research and Interdisciplinary Programs with the School of Engineering, Santa Clara University, Santa Clara, CA, USA.
Dr. Kitts is a fellow of the American Society of Mechanical Engineers. \end{IEEEbiography}

\begin{thebibliography}{99}

\bibitem{1} J. Cortés and M. Egerstedt, "Coordinated control of multi-robot systems: A survey," \textit{SICE Journal of Control, Measurement, and System Integration}, vol. 10, no. 6, pp. 495--503, 2017.

\bibitem{2} J. McDonald, C. A. Kitts, and M. A. Neumann, "Experimental implementation and verification of scalar field ridge, trench, and saddle point maneuvers using multirobot adaptive navigation," \textit{IEEE Systems Journal}, vol. 14, no. 2, pp. 2668--2679, 2019.

\bibitem{3} F. Pagano, A. Renzaglia, L. Brinon-Arranz, and L. Schenato, "Distributed multi-robot active-sensing of a diffusive source," \textit{IEEE Robotics and Automation Letters}, 2025.

\bibitem{4} P. Ögren, E. Fiorelli, and N. E. Leonard, "Cooperative control of mobile sensor networks: Adaptive gradient climbing in a distributed environment," \textit{IEEE Transactions on Automatic Control}, vol. 49, no. 8, pp. 1292--1302, 2004.

\bibitem{5} S. Li, R. Kong, and Y. Guo, "Cooperative distributed source seeking by multiple robots: Algorithms and experiments," \textit{IEEE/ASME Transactions on Mechatronics}, vol. 19, no. 6, pp. 1810--1820, 2014.

\bibitem{6} Z. Jin, Y. Hu, and X. Zhang, "Gradient-free cooperative source-seeking of quadrotor under disturbances and communication constraints," \textit{IEEE Transactions on Industrial Electronics}, 2024.

\bibitem{7} M. Krstic and H.-H. Wang, "Stability of extremum seeking feedback for general nonlinear dynamic systems," \textit{Automatica}, vol. 36, no. 4, pp. 595--601, 2000.

\bibitem{8} C. Cooper, C. Waight, and C. A. Kitts, "Initial study of multirobot adaptive navigation for vector field mapping," in \textit{Proc. ASME International Mechanical Engineering Congress and Exposition}, 2019.

\bibitem{9} D. Panagou, "Motion planning and collision avoidance using navigation vector fields," in \textit{Proc. IEEE International Conference on Robotics and Automation}, 2014, pp. 2513--2518.

\bibitem{10} F. Gul, I. Mir, W. Rahiman, and S. M. Islam, "A consolidated review of path planning and optimization techniques: Technical perspectives and future directions," \textit{Electronics}, vol. 10, no. 18, p. 2250, 2021.

\bibitem{11} L. Briñón-Arranz, A. Renzaglia, and L. Schenato, "Multirobot symmetric formations for gradient and hessian estimation with application to source seeking," \textit{IEEE Transactions on Robotics}, vol. 35, no. 3, pp. 782--789, 2019.

\bibitem{12} Z. Deng and J. Luo, "Fully distributed algorithms for constrained nonsmooth optimization problems of general linear multiagent systems and their application," \textit{IEEE Transactions on Automatic Control}, vol. 69, no. 2, pp. 1377--1384, 2023.

\bibitem{13} S. T. Hart, J. Kamenetsky, and C. A. Kitts, "Dynamic elliptical shaping control for swarm robots," \textit{IEEE Access}, vol. 11, pp. 17454--17470, 2023.

\bibitem{14} C. A. Kitts and I. Mas, "Cluster space specification and control of mobile multirobot systems," \textit{IEEE/ASME Transactions on Mechatronics}, vol. 14, no. 2, pp. 207--218, 2009.

\bibitem{15} T. Adamek, C. A. Kitts, and I. Mas, "Gradient-based cluster space navigation for autonomous surface vessels," \textit{IEEE/ASME Transactions on Mechatronics}, vol. 20, no. 2, pp. 506--518, 2015.

\bibitem{16} I. Mas, S. Li, J. Acain, and C. A. Kitts, "Entrapment/escorting and patrolling missions in multi-robot cluster space control," in \textit{Proc. IEEE/RSJ International Conference on Intelligent Robots and Systems}, 2009, pp. 5855--5861.

\bibitem{17} S. T. Hart and C. A. Kitts, "Unifying control architecture for reactive particle swarms," \textit{IEEE/ASME Transactions on Mechatronics}, vol. 28, no. 2, pp. 873--883, 2022.

\bibitem{18} S. T. Hart, J. Kamenetsky, C. A. Kitts, and I. Mas, "Circular formation control of multiple unicycle-class robots," \textit{IEEE Transactions on Control Systems Technology}, vol. 26, no. 3, pp. 1140--1147, 2018.

\bibitem{19} H. Cao, S. Shreedharan, and N. Atanasov, "Multi-robot object SLAM using distributed variational inference," \textit{IEEE Robotics and Automation Letters}, vol. 9, no. 3, pp. 2763--2770, 2024.

\bibitem{20} P. S. Trakas, A. Tantoulas, and C. P. Bechlioulis, "Formation control of nonlinear multi-agent systems with nested input saturation," \textit{Applied Sciences}, vol. 14, no. 1, p. 213, 2023.

\bibitem{21} C. Waight and C. A. Kitts, "A functional test bed for multi-robot adaptive navigation in vector fields," in \textit{Proc. ASME International Design Engineering Technical Conferences and Computers and Information in Engineering Conference}, 2025.

\bibitem{22} R. K. Lee, C. A. Kitts, and M. A. Neumann, "Multiple UAV adaptive navigation for three-dimensional scalar fields," \textit{IEEE Access}, vol. 9, pp. 122626--122654, 2021.

\bibitem{23} R. K. Lee, C. A. Kitts, and M. A. Neumann, "Structured isosurface mapping of 3-D scalar fields while seeking extremum with mobile sensor networks," \textit{IEEE Sensors Journal}, vol. 24, no. 6, pp. 7325--7335, 2024.

\bibitem{24} S. Tomer, D. J. Edwards, and C. A. Kitts, "A low-cost indoor testbed for multirobot adaptive navigation research," in \textit{Proc. IEEE/ASME International Conference on Advanced Intelligent Mechatronics}, 2014, pp. 140--145.

\end{thebibliography}




\end{document}
